\section{Method details}
\label{app:methods}

This appendix records the scheme checks, implementation contracts, and routing details underlying Section~\ref{sec:methods}.

\subsection{Certified method mixing}

Mixing methods offers three extensions of increasing aggressiveness. The
reported experiments do not rely on them, and the soundness theorem for the
base grammar does not automatically cover them.

\textbf{Averaged blending.} A convex combination
$\mathrm{avg}(\lambda; T_1, T_2) = \lambda T_1 + (1-\lambda) T_2$ of two
averaged operators is averaged with a computable constant
\citep[\S2.3]{RyuYin2022LargeScale}, given in closed form by
\citet{CombettesYamada2015Compositions}.
Lemma~\ref{lem:blend} in Appendix~\ref{app:theory} gives a certificate rule
when both maps share a state and metric and have a nonempty common
fixed-point set. A common fixed point is a semantic obligation not established
by the present syntactic scheme gate. Thus the lemma describes a conditional
extension, not an automatically admitted genome covered by
Theorem~\ref{thm:soundness}.

\textbf{Three-operator splitting.} Davis--Yin splitting
\citep{DavisYin2017Three} is the canonical certified fusion of
forward--backward and Douglas--Rachford, handling all three blocks of
problem~\eqref{eq:composite} simultaneously; it enters the genome as one
additional base scheme with its own gated side condition.

\textbf{Safeguarded switching.} A $\mathrm{switch}(\text{trigger},
\text{aggressive}, \text{fallback})$ gene runs an aggressive component
(stepsizes outside the certified region, Anderson-type extrapolation,
aggressive restarts) under a summable-error safeguard. Relative to the
certified fallback step, the accepted perturbation at iteration $k$ must have
metric norm at most a prescribed $\eta_k$ with $\sum_k\eta_k<\infty$;
otherwise the fallback step is taken. The resulting inexact iteration
therefore inherits the fallback's convergence (Section~\ref{sec:theory}). The hand-designed
literature illustrates distinct ways to compose solver decisions: HPR-LP
uses Halpern-accelerated semi-proximal Peaceman--Rachford splitting
\citep{Chen2024HPRLP}, whereas PDLP adapts the primal weight of a PDHG-based
solver \citep{Applegate2021Practical}. A switch gene would make guarded
versions of such decisions searchable; its convergence claim depends on the
stated summable-error condition.

Blends and safeguarded switching are promising directions for discovering
compositions beyond single schemes. They are not exercised in
Section~\ref{sec:experiments}; blending requires a common-fixed-point
witness or a restricted constructive typing rule before it can be admitted
with the guarantee claimed for the formal grammar.


\subsection{The Stage-0 type-check gate}

For the base-scheme and relaxation grammar of Section~\ref{sec:theory},
the Stage-0 gate checks the convergence side condition before execution:
$\alpha \in (0, 2/L_f)$ for FBS and Davis--Yin; any $\alpha > 0$
for DRS;
$\tau\sigma\ell^2 < 1$ for PDHG;
$\tau L_f/2 + \tau\sigma\ell^2 < 1$ for Condat--V\~u, where
$\ell \geq \|L\|$ is the supplied operator-norm bound
\citep[\S2.7 and \S3.3]{RyuYin2022LargeScale}; and for the wrappers,
$\rho\,\theta < 1$ for over-relaxation of a
$\theta$-averaged map, with Halpern anchoring admitted only
for nonexpansive iterations. FBS, DRS, and Davis--Yin additionally require
the exact proximal maps of their chosen splitting to be available
(Appendix~\ref{app:theory}). A genome that satisfies its conditions receives
a construction certificate recording the checked inequality, the
resulting averagedness constant, and the textbook citation; a genome that
does not is rejected, never silently clamped, and the violated
inequality is returned verbatim as feedback to the proposer. The scheme gate
is re-run inside every solve; a 20{,}000-genome fuzz found zero disagreements
between that gate and the solver-level checks. The theorem also assumes an
admissible formulation and exact atoms: Section~\ref{sec:theory}
distinguishes these mathematical obligations from numerical contract tests.


\subsection{Compilation and certificates}

An admitted scheme is compiled to a jitted JAX iteration; the mathematical
certificate applies to its exact reference update, not automatically to
floating-point execution. Genome search never executes LLM-written numerical
code. LLM-written code runs only in operator synthesis
(Section~\ref{sec:synthesis}), where a proposed projection is executed only
inside its contract harness and enters a solve only after passing it. Beyond the construction certificate
(level 1), each solve carries runtime certificates (level 3): the step residual
$\|z_{k+1}-z_k\|_{H_G}$ of Section~\ref{sec:theory} and, for
problem~\eqref{eq:composite} with strongly convex $f$, a rigorous duality
gap. For TV denoising the dual value at any
$u$ with $\|u\|_\infty \le \lambda$ is
$\langle y, D^{\!\top}u\rangle - \tfrac12\|D^{\!\top}u\|^2$, so
$\mathrm{gap}(x, u) \ge 0$ always and termination at relative gap
$\le \epsilon$ is a checkable certificate of $\epsilon$-optimality, not a
heuristic stopping rule. The benchmark record also stores the original-problem acceptance result.


\subsection{Hardware-aware backend genes}

A second family of genes controls not what iteration runs but
how it is computed: floating-point precision, batching across
instances, sparse-matrix format, and kernel-fusion strategy. The construction
certificate concerns an exact operator $T$, not the backend's floating-point
behavior. Fusion, precision, and block sizes can change numerical results,
so the implementation is checked separately by
per-atom tests against a reference backend to tolerance, in the pattern of
hardware-aware kernel generation for coding agents
\citep{Hawkeye2026}. This separates two implementation-level checks from
the certified core: a safeguard monitor for schedules above the calculus and
unit tests for kernels below it. Neither check alone proves the
exact-operator assumptions of Section~\ref{sec:theory}.

The target hardware is visible to the search: the device identifier
conditions the mutation context, so genomes can specialize to the
accelerator they will run on. Kernel fusion matters most exactly where
per-iteration work is small and launch overhead dominates, i.e.\ the
microsecond-budget cells; batched execution matters where many instances
are solved simultaneously. This paper ships the precision, batching, and
fused-versus-naive step genes; autogenerated fused kernels for the atoms are
future work.

\paragraph{The fusion tile is a gene.}
A temporally fused kernel runs $T$ iterations per launch on a $(B+2T) \times
(B+2T)$ tile whose halo is held in shared memory, so the pair $(B, T)$ is a
backend gene like precision or batching. Three things make it a gene rather
than a constant. Its admissible set is device dependent: a tile whose halo
exceeds the device's opt-in shared memory per block cannot launch at all, which
removes a third of the grid on a consumer part and none of it on a
data-center part. Its correctness is a contract, not an assumption: an
admitted tile must reproduce $T$ sequential unfused iterations, which at
relaxation $\rho = 1$ means bitwise, so the gene cannot trade accuracy for
speed. And its value is measured, not derived: the cost model prices a tile by
the per-iteration cost it achieves on the device in question, exactly as it
prices an operator, so the tile competes inside
equation~\eqref{eq:costmodel} rather than being fixed ahead of the search. At
solve time the device is identified from its reported properties (name, compute
capability, multiprocessor count, shared memory per block) and the tile is
selected for that device; declining to fuse, that is, running the compiled loop
of the array compiler, is always one of the candidates.

\paragraph{The device is a candidate.}
Because an exact operator leaves the iteration count $N$ unchanged, the same
(formulation, operator) pair on a machine's GPU and on its CPU are two
candidates that share one measured $N$ and differ only in setup and per-iteration
cost. The router therefore enumerates (formulation, operator, device) triples
and prices each device from its own calibration; a machine's CPU is a separate
device from its GPU, and a device for which the process has no backend is
declined rather than priced from a borrowed measurement.
Section~\ref{sec:exp:device} reports where the device, the tile and the
precision decision change the winner.


\subsection{Formulation genes}

The same feasible set admits many splittings of problem~\eqref{eq:composite},
and which constraints live in the linear operator $L$ and which live in the
proximal step is itself a gene. Relocation removes a constraint family from
$h(Lx)$ and enforces it exactly through the projection in $g$: a single dense
row together with the variable bounds it touches (a box intersected with a
hyperplane or a slab), a product of disjoint such blocks, or a family of
second-order cones. The kept rows stay in $L$ and are optionally equilibrated
row by row. A planner reads the candidate families off $(P, q, A, l, u)$ and
works under any row and column permutation of the instance, so nothing depends
on how a benchmark happens to order its variables. Relocation changes the
iteration but not the problem, and the certificate transfers only when the
projection is exact; the planner declines when no exact projection exists.
The formulations the router can currently choose are the standard form with
and without Ruiz equilibration, row relocation with and without row scaling,
product relocation under several row-selection rules, and cone relocation.
Each transformation must preserve the original solution set under its recovery
map and pass the scheme gate with its own operator bound; the formal transfer
rule is Proposition~\ref{prop:formulation-transfer}. Tests of an implemented
projection are a separate obligation from its mathematical exactness.


\subsection{Operator synthesis under contract}

A relocated formulation is only as fast as its projection, and when the cost
model finds that the projection's per-call time is the binding term it emits a
specification rather than a guess: the set (for example
$\{x : lo \le x \le hi,\ a^{\top}x \in [b_{lo}, b_{hi}]\}$), the device, the
per-call budget that would change the routing decision, and a contract of
input conditions the implementation must handle (mixed-sign $a$, infinite
bounds, degenerate breakpoint ties, very small and very large dimension, and
the inputs seen in deployment). An LLM evolves implementations against that
specification inside SkyDiscover \citep{Liu2026SkyDiscover}. A candidate
implementation is admitted only if it passes the contract tests: agreement
with a reference projection on every tested clause. These tests
screen numerical code but do not prove that it equals the exact projection
on all inputs. Admitted operators are calibrated per device and stored with
their provenance. A projection that is fast but fails the contract tests is
never used, however good its score.


\subsection{The gene pool, the cost model, and routing}

Evolved components, formulations, measurements, and lessons are kept in an
append-only store, the gene pool. The router decomposes the time to a verified
answer as
\begin{equation}
  T = \text{setup} + N(\text{formulation}) \cdot t(\text{operator}, \text{device}),
  \label{eq:costmodel}
\end{equation}
where the iteration count $N$ is taken from measured evidence for the
formulation and never from an operator-norm bound, and the per-iteration cost
$t$ and the setup are taken from calibration of the operator on the device.
Candidates are the product of formulations, contract-tested projection
implementations, and devices, together with external engines, each run in its
own interpreter: interior-point
and first-order generic solvers \citep{Goulart2024Clarabel,Schwan2023PIQP,Stellato2020OSQP,ODonoghue2016Conic,Domahidi2013ECOS,HiGHS2026Solvers,Applegate2021Practical},
and structure-specific solvers that are offered only after a detector has
recognized their structure in the instance under permutation: coordinate
descent for lasso and group lasso \citep{Massias2018Celer,Bertrand2022Beyond,Pedregosa2011Scikit},
network simplex and minimum-cost flow for network linear programs
\citep{Bonneel2011Displacement,Flamary2021POT}, and a Frank--Wolfe core-set
method for the minimum enclosing ball \citep{Yildirim2008Two,LacosteJulien2015Global}.
A candidate without evidence is explored, and its run is charged to the budget;
exploration stops once it has cost twice the time to the first verified answer,
and the verified incumbent is never discarded. Every answer, whichever engine
produced it, is accepted only if it passes the same check on the original
problem. An LLM agent drives the system through four actions: solve an
instance, evolve an operator against a specification, propose a planner rule,
and propose a cost-model term, each admitted only after measurement on held-out
instances.


\subsection{Evolutionary loop}

Search proceeds over generations with elitism. Two proposal engines operate
on the identical genome space at identical evaluation budgets: (i) an
LLM backend that receives the current archive (best genomes with their
metrics), the mutation menu, and the genome JSON schema, and replies with a
patched genome; invalid JSON or a Stage-0 rejection triggers one retry
with the error injected, then a fallback to the random proposer; proposals
are deterministically seeded and the archive deduplicates by content hash;
and (ii) a random-search control that rejection-samples the valid
region with log-uniform stepsizes. Fitness is the negated shifted geometric
mean (SGM-10) of wall-clock time-to-certificate over the training instances;
runs that fail to reach the certificate within the iteration budget score
$-\infty$. The loop monitors for degenerate (flat) fitness landscapes and
reports them rather than proceeding silently. Final comparisons are made
exclusively on held-out instances never seen during search. The backend
interface is engine-agnostic by construction, so the evolutionary machinery
itself (island scheduling, parent sampling) can be swapped without touching
the gate, compiler, or certificates \citep{Liu2026SkyDiscover}.


\subsection{Implementation}

Solvers are compiled with JAX. The total-variation search of
Section~\ref{sec:exp:evolve} runs at float64 on the CPU of a single
workstation, with Qwen2.5-Coder-1.5B-Instruct served by vLLM on its RTX~4090 as
the mutation model. Operator synthesis uses Qwen3.8-27B in FP8 served by vLLM
on an H100. Solves are timed on an RTX~4090 workstation and on an H100 80\,GB
node; baselines and specialists run on the CPU of the machine being measured,
in separate interpreters that never share an environment with our solver.

